\section*{Acknowledgements}

This work was developed in the context of the PLURALISMO project PLU230018, \emph{Evaluación de Estrategias para Diseminar Fake News Usando Inteligencia Artificial}. The authors thank the contributors to the earlier endorsement-based simulation studies on which the reasoning architecture builds.

\section*{Disclosure statement}

The authors report no competing interests. This statement must be confirmed by all authors before submission.

\section*{Funding}

Project identifier PLU230018 is documented in the software repository. The formal funder name, award wording, and any required grant acknowledgement must be confirmed before journal submission.

\section*{Data and code availability}

The FAKENEWS-ABM source code and experiment harness are available at \url{https://github.com/pleger/FAKENEWS-ABM}. The analyses reported here use these experiment roots:
\begin{itemize}
  \item \path{research_questions_20260812_180619};
  \item \path{dissemination_experiments_20260813_215632}; and
  \item \path{dissemination_experiments_20260814_083018}.
\end{itemize}
Before submission, the exact code commit and generated analysis tables should be archived in a versioned release or research repository with a persistent identifier.

\section*{Author contributions}

The final CRediT contribution statement requires approval from all authors. Proposed roles to verify are: conceptualisation, methodology, software, validation, formal analysis, investigation, data curation, writing---original draft, writing---review and editing, visualisation, supervision, project administration, and funding acquisition. No individual allocation is asserted in this draft without author confirmation.

\section*{Ethics statement}

The study analyses synthetic agents and simulation outputs and does not use human participants or personal data. Any future empirical calibration dataset must undergo its own ethical and privacy review.
